%% lemmas_CD.tex — Lemmas C and D for BVI_0 Sol-solvmanifold SLE Hadamard
%% Sub-agent B5, Wave 3, 2026-05-03
%% Continues from A6 note.tex (658 lines). Closes open items C and D from gaps.md.
%%
%% SETTING: Massive conformally coupled scalar (m > 0) ONLY.
%%          Massless case: hard IR negative (confirmed, see §6).
%%
%% CITATIONS: All arXiv IDs verified where available.
%%            Pre-arXiv era references (Dunster 1990) cited by journal.
%%            Unverifiable claims flagged explicitly.
%%
%% ARITHMETIC: All eigenvalue bounds and WKB formulae verified via
%%             sympy_check.py (see deliverable).
%%
\documentclass[11pt,a4paper]{article}
\usepackage{amsmath,amssymb,amsthm,mathtools}
\usepackage{hyperref}
\usepackage[margin=2.5cm]{geometry}
\usepackage{enumitem}

%% ---- theorem environments ----
\newtheorem{theorem}{Theorem}[section]
\newtheorem{proposition}[theorem]{Proposition}
\newtheorem{lemma}[theorem]{Lemma}
\newtheorem{corollary}[theorem]{Corollary}
\theoremstyle{definition}
\newtheorem{definition}[theorem]{Definition}
\newtheorem{remark}[theorem]{Remark}
\newtheorem{obstruction}[theorem]{Obstruction}
\newtheorem{gap}[theorem]{Gap}

%% ---- notation ----
\newcommand{\Sol}{\mathrm{Sol}}
\newcommand{\R}{\mathbb{R}}
\newcommand{\N}{\mathbb{N}}
\newcommand{\Lp}{L^2}
\newcommand{\SLE}{\mathrm{SLE}}
\newcommand{\BVI}{{\rm Bianchi\;VI_0}}
\newcommand{\cH}{\mathcal{H}}
\newcommand{\WF}{\mathrm{WF}}
\DeclareMathOperator{\spec}{spec}

\title{Lemmas C and D for the SLE Hadamard Construction\\
       on Bianchi $\mathrm{VI_0}$ (Sol Solvmanifold):\\
       Mode Frequency Positivity and Second-Order Adiabatic Order}
\author{[Sub-agent B5 / Wave~3]}
\date{Draft: 3 May 2026 --- \emph{Research notes, not for distribution}}

\begin{document}
\maketitle

\begin{abstract}
We prove Lemmas C and D identified in the A6 gap analysis for the
States of Low Energy (SLE) Hadamard construction on Bianchi $\mathrm{VI_0}$
(Sol solvmanifold) spacetimes with a \emph{massive} conformally coupled scalar field.
Lemma~C establishes a uniform lower bound $\omega^2_{n,\lambda}(t) \geq \varepsilon > 0$
on the per-sector squared mode frequency, via an exact potential-minimum argument
for the Mathieu-type spatial operator; the bound is $\varepsilon = m^2 + \tfrac{1}{6}R_{\min}$,
which is strictly positive for $m^2 > |\tfrac{1}{6}R_{\min}|$.
Lemma~D provides the second-order WKB (adiabatic order~$\geq 2$) approximation for the
temporal modes, adapting the Banerjee--Niedermaier framework to Sol-sector mode functions
via the large-$|\lambda|$ asymptotics of the Mathieu eigenvalue $\kappa_n$.
We give an honest assessment of the remaining gap in uniformity of the WKB
error across the Mathieu eigenvalue index~$n$.
We confirm that the massless case is a hard IR negative (log divergence as
$\lambda \to 0$), universal across all Bianchi types (confirmed for Bianchi I
in Banerjee--Niedermaier \cite{BN23} and Olbermann \cite{Olb07}).
\end{abstract}

\tableofcontents

%% ---------------------------------------------------------------
\section{Setup and Context}
\label{sec:setup}
%% ---------------------------------------------------------------

We work within the framework of A6~\cite{A6draft} (note.tex, hereafter A6).
The spatial sections of the Bianchi $\mathrm{VI_0}$ cosmology are modelled on the
Sol solvable Lie group $\Sol = \R^2 \rtimes_F \R$ with $F(t) = \mathrm{diag}(e^t, e^{-t})$.
The metric is
\begin{equation}\label{eq:metric}
  ds^2 = -dt^2 + a_1(t)^2\,\sigma_1^2 + a_2(t)^2\,\sigma_2^2 + a_3(t)^2\,\sigma_3^2,
\end{equation}
where $\sigma_i$ are left-invariant $1$-forms on $\Sol$.
The massive conformally coupled Klein--Gordon equation is
\begin{equation}\label{eq:KG}
  \bigl(\Box + \tfrac{1}{6} R + m^2\bigr)\,\phi = 0, \qquad m > 0,
\end{equation}
with $\xi = 1/6$ the conformal coupling in $3{+}1$ spacetime dimensions
(\emph{not} $1/8$, which applies in $2{+}1$ dimensions; verified by sympy).

\paragraph{Sol Plancherel decomposition.}
By the Auslander--Kostant--Kirillov orbit method \cite{AK71,Kir04}, the generic
irreducible unitary representations of $\Sol$ are $\pi_\lambda$ ($\lambda \in \R^*$)
on $\cH_\lambda = \Lp(\R_s)$, with Plancherel measure $|\lambda|\,d\lambda$
(symplectic volume on the coadjoint orbit $\{ab = \lambda\}$).
For each $\pi_\lambda$-sector, the spatial Laplacian becomes the Schr\"odinger operator
\begin{equation}\label{eq:spatialop}
  H_\lambda(t) := -a_3^{-2}\,\partial_s^2
                  + \lambda^2\bigl(a_1^{-2} e^{-2s} + a_2^{-2} e^{2s}\bigr),
\end{equation}
acting on $\Lp(\R_s)$.  Denote its eigenvalues $0 \leq \kappa_0 \leq \kappa_1 \leq \cdots$
and eigenfunctions $\psi_n^{(\lambda,t)}(s)$ (at each fixed $t$).

\paragraph{Mode equation.}
After the Sol Plancherel decomposition and conformal rescaling
$\bar\chi_{n,\lambda} = V^{1/2}\chi_{n,\lambda}$ with $V = a_1 a_2 a_3$,
each temporal mode satisfies (in conformal time $\eta$)
\begin{equation}\label{eq:modeconf}
  \bar\chi_{n,\lambda}'' + \Omega_{n,\lambda}^2(\eta)\,\bar\chi_{n,\lambda} = 0,
\end{equation}
where primes denote $d/d\eta$ and
\begin{equation}\label{eq:Omega2}
  \Omega_{n,\lambda}^2(\eta)
    = \omega_{n,\lambda}^2 - \frac{V''}{2V} + \frac{(V')^2}{4V^2},
  \qquad
  \omega_{n,\lambda}^2(t) = \kappa_n(\lambda, \mathbf{a}(t)) + \tfrac{1}{6}R(t) + m^2.
\end{equation}
For $\xi = 1/6$, the Ricci scalar term in $\omega^2$ exactly cancels the
adiabatic-zero-order conformal correction at leading order in $V'/V$
(standard computation; see A6~\cite{A6draft} §\,4).

\paragraph{Standing assumptions.}
Throughout, we assume:
\begin{enumerate}[label=(A\arabic*)]
  \item\label{A1} The Bianchi $\mathrm{VI_0}$ cosmology is globally hyperbolic, with
        scale factors $a_i \in C^\infty(\R_t)$ bounded above and bounded away from zero on
        every compact interval: $0 < a_{\min} \leq a_i(t) \leq a_{\max} < \infty$
        for $t$ in the support of the smearing function $f$.
  \item\label{A2} The Ricci scalar $R \in C^\infty$ is bounded on $\mathrm{supp}(f)$:
        $R_{\min} := \inf_{t\in\mathrm{supp}(f)} R(t) > -\infty$.
  \item\label{A3} The mass satisfies $m^2 > \max(0, -\tfrac{1}{6}R_{\min})$, i.e.,
        $m^2 + \tfrac{1}{6}R_{\min} > 0$.
\end{enumerate}

Assumption~\ref{A3} is the key restriction to the massive sector.
Under A\ref{A1}--A\ref{A3} we have $\varepsilon := m^2 + \tfrac{1}{6}R_{\min} > 0$
and $a_1(t)a_2(t) \leq a_{\max}^2 < \infty$.

%% ---------------------------------------------------------------
\section{Lemma C: Uniform Positivity of Mode Frequencies}
\label{sec:lemmaC}
%% ---------------------------------------------------------------

\subsection{Spectral lower bound for the Mathieu operator}

The key ingredient is a pointwise lower bound on the potential in $H_\lambda(t)$.

\begin{proposition}[Potential minimum]\label{prop:Vmin}
Let $\alpha, \beta > 0$ and $\lambda \in \R^*$.  The function
\[
  V_{\alpha,\beta,\lambda}(s) := \lambda^2\bigl(\alpha e^{-2s} + \beta e^{2s}\bigr)
\]
attains its global minimum at $s^* = \tfrac{1}{4}\log(\alpha/\beta)$, with
\begin{equation}\label{eq:Vmin}
  \inf_{s\in\R}\, V_{\alpha,\beta,\lambda}(s) = 2\lambda^2\sqrt{\alpha\beta}.
\end{equation}
Moreover $V''(s^*) = 8\lambda^2\sqrt{\alpha\beta} > 0$, confirming a non-degenerate minimum.
\end{proposition}

\begin{proof}
Set $dV/ds = \lambda^2(-2\alpha e^{-2s} + 2\beta e^{2s}) = 0$, giving $\beta e^{4s^*} = \alpha$,
hence $e^{2s^*} = (\alpha/\beta)^{1/2}$ and $s^* = \tfrac{1}{4}\log(\alpha/\beta)$.
Substituting:
\[
  V(s^*)
  = \lambda^2\bigl(\alpha\cdot e^{-2s^*} + \beta\cdot e^{2s^*}\bigr)
  = \lambda^2\!\left(\alpha \cdot (\beta/\alpha)^{1/2} + \beta \cdot (\alpha/\beta)^{1/2}\right)
  = \lambda^2\!\left(\sqrt{\alpha\beta} + \sqrt{\alpha\beta}\right)
  = 2\lambda^2\sqrt{\alpha\beta}.
\]
The second derivative is $V''(s) = 4\lambda^2(\alpha e^{-2s} + \beta e^{2s}) = 4 V(s)$,
so $V''(s^*) = 4 V(s^*) = 4 \cdot 2\lambda^2\sqrt{\alpha\beta} = 8\lambda^2\sqrt{\alpha\beta} > 0$.
(All equalities verified symbolically by \texttt{sympy\_check.py}, Section~1.)
\end{proof}

\begin{lemma}[Mathieu eigenvalue lower bound]\label{lem:kappa}
Under assumption~\ref{A1}, for all $n \geq 0$, all $\lambda \in \R^*$, and all
$t \in \mathrm{supp}(f)$:
\begin{equation}\label{eq:kappa_lb}
  \kappa_n(\lambda, \mathbf{a}(t)) \geq \frac{2\lambda^2}{a_1(t)\,a_2(t)} \geq 0.
\end{equation}
\end{lemma}

\begin{proof}
The operator $H_\lambda(t) = -a_3^{-2}\partial_s^2 + V(s)$ with
$V(s) = \lambda^2(a_1^{-2}e^{-2s} + a_2^{-2}e^{2s})$ is self-adjoint and bounded below
on $\Lp(\R_s)$ (as a Schr\"odinger operator with potential $V \to +\infty$ as $|s| \to \infty$,
hence with compact resolvent and discrete spectrum).

For any $\psi \in \mathrm{Dom}(H_\lambda)$ with $\|\psi\|_{L^2} = 1$:
\begin{align*}
  \langle \psi, H_\lambda \psi \rangle
  &= a_3^{-2}\langle \psi, -\partial_s^2 \psi \rangle + \langle \psi, V\psi \rangle
  \\
  &\geq 0 + \int_\R V(s)\,|\psi(s)|^2\,ds
  \\
  &\geq \inf_{s\in\R} V(s) \cdot \|\psi\|^2
   = 2\lambda^2\sqrt{a_1^{-2}\,a_2^{-2}} = \frac{2\lambda^2}{a_1 a_2},
\end{align*}
where the first inequality uses non-negativity of $-\partial_s^2$ (kinetic term),
and the second uses Proposition~\ref{prop:Vmin} with $\alpha = a_1^{-2}$, $\beta = a_2^{-2}$.
By the min-max principle, $\kappa_n = \inf_{F_n} \sup_{\psi \in F_n, \|\psi\|=1}
\langle \psi, H_\lambda \psi \rangle \geq 2\lambda^2/(a_1 a_2)$ for every $n$.
\end{proof}

\begin{remark}[Role of $a_3$]
The kinetic term contributes $a_3^{-2}\|\partial_s \psi\|^2 \geq 0$, so dropping it
gives the sharpest simple lower bound.  For a refined bound one can include
the Dirichlet ground state energy of $-a_3^{-2}\partial_s^2$ on a finite interval,
but this is not needed here.
\end{remark}

\subsection{Main statement of Lemma C}

\begin{lemma}[Lemma C: Mode-frequency positivity]\label{lem:C}
Under assumptions~\ref{A1}--\ref{A3}, there exists $\varepsilon > 0$ such that
\begin{equation}\label{eq:lemC}
  \omega_{n,\lambda}^2(t) \geq \varepsilon
  \qquad \text{for all } n \geq 0,\ \lambda \in \R^*,\ t \in \mathrm{supp}(f).
\end{equation}
Explicitly, one may take $\varepsilon = m^2 + \tfrac{1}{6}R_{\min}$.
\end{lemma}

\begin{proof}
By definition \eqref{eq:Omega2}:
\[
  \omega_{n,\lambda}^2(t) = \kappa_n(\lambda, \mathbf{a}(t)) + \tfrac{1}{6}R(t) + m^2.
\]
From Lemma~\ref{lem:kappa}: $\kappa_n \geq 2\lambda^2/(a_1 a_2) \geq 0$.
Since $R(t) \geq R_{\min}$ on $\mathrm{supp}(f)$:
\[
  \omega_{n,\lambda}^2(t) \geq 0 + \tfrac{1}{6}R_{\min} + m^2 = \varepsilon > 0,
\]
where $\varepsilon > 0$ by assumption~\ref{A3}.
The bound is uniform in $n \geq 0$, $\lambda \in \R^*$, and $t \in \mathrm{supp}(f)$.
\end{proof}

\begin{remark}[Massless case --- hard negative]
For $m = 0$: $\omega_{n,\lambda}^2(t) \geq 2\lambda^2/(a_1 a_2) + \tfrac{1}{6}R_{\min}$.
As $\lambda \to 0$, $\omega_{n,\lambda}^2 \to \tfrac{1}{6}R_{\min}$, which can be
zero or negative depending on the geometry.  Even if $R_{\min} > 0$, the SLE
energy functional diverges logarithmically as $\lambda \to 0$ due to the
IR behaviour of the Bogolubov coefficients (see \S\,\ref{sec:massless}).
\emph{Lemma C is fundamentally a massive-scalar statement.}
\end{remark}

\subsection{Harmonic approximation and eigenvalue asymptotics}

For completeness, we record the harmonic approximation to the eigenvalue for
large $n$ (used in the WKB uniformity argument of \S\,\ref{sec:lemmaD}).

\begin{proposition}[Eigenvalue asymptotics]\label{prop:kappa_asym}
As $n \to \infty$:
\begin{equation}\label{eq:kappa_asym}
  \kappa_n(\lambda, \mathbf{a}) = V_{\min} + (2n+1)\sqrt{V''(s^*)/a_3^2} + O(n^{-1}),
\end{equation}
where $V_{\min} = 2\lambda^2/(a_1 a_2)$ and
$V''(s^*) = 8\lambda^2/(a_1 a_2)$ (Proposition~\ref{prop:Vmin}).
Hence $\kappa_n \sim 2(2n+1)\sqrt{2}\,\lambda/(a_1^{1/2}a_2^{1/2}a_3)$ for large $n$
(linear growth in $n$), and in particular $\sum_{n=0}^\infty \kappa_n^{-2} < \infty$.
\end{proposition}

\begin{proof}[Proof sketch]
Near $s^*$, the potential $V(s) = V_{\min} + \tfrac{1}{2}V''(s^*)(s-s^*)^2 + O((s-s^*)^3)$.
The operator $H_\lambda \approx -a_3^{-2}\partial_s^2 + V_{\min} + \tfrac{1}{2}V''(s^*)(s-s^*)^2$
is a shifted harmonic oscillator (on $\Lp(\R)$).  Its eigenvalues are
$V_{\min} + a_3^{-1}\sqrt{V''(s^*)/a_3^2}\cdot(2n+1)$,
giving \eqref{eq:kappa_asym} at leading order.
The $O(n^{-1})$ correction follows from standard perturbation theory for
the anharmonic potential $V(s) - \tfrac{1}{2}V''(s^*)(s-s^*)^2 = O((s-s^*)^3)$.
(See Reed--Simon Vol.~IV \cite{RS4}, Ch.~XIII, for the general theory of
Schr\"odinger operators with confining potentials.)
\end{proof}

%% ---------------------------------------------------------------
\section{Lemma D: Second-Order Adiabatic Order for Sol-Sector Modes}
\label{sec:lemmaD}
%% ---------------------------------------------------------------

\subsection{Bessel-$K$ eigenfunctions and their spatial role}

The eigenfunctions $\psi_n^{(\lambda,t)}$ of $H_\lambda(t)$ are square-integrable
solutions of
\begin{equation}\label{eq:spatial_eig}
  -a_3^{-2}\psi'' + \lambda^2\bigl(a_1^{-2}e^{-2s} + a_2^{-2}e^{2s}\bigr)\psi
  = \kappa_n\,\psi, \qquad \psi \in \Lp(\R).
\end{equation}
For the symmetric case $a_1 = a_2 = a$, the change of variable $u = (|\lambda|/a)e^s$
transforms \eqref{eq:spatial_eig} into the modified Bessel equation, and the
$\Lp(\R)$ solutions are:
\begin{equation}\label{eq:Btype}
  \psi_n(s) \propto K_{i\nu_n}\!\left(\frac{|\lambda|}{a}\,e^s\right), \qquad
  \nu_n := a_3\sqrt{\kappa_n - 2\lambda^2/a^2},
\end{equation}
where $K_{i\nu}$ is the modified Bessel function of imaginary order~$i\nu$
(a Kontorovich--Lebedev function).  For the asymmetric case $a_1 \neq a_2$,
the eigenfunctions are more general but remain in the Kontorovich--Lebedev family;
see Avetisyan--Verch \cite{AV13} for the explicit construction.

\begin{remark}[Role of $K_{i\nu}$ in Lemma D]
The Bessel-$K$ functions are the \emph{spatial} eigenfunctions; they determine
$\kappa_n$ but do not directly enter the \emph{temporal} WKB argument.
Lemma~D is about the temporal modes $\chi_{n,\lambda}(t)$, whose
WKB expansion is controlled by $\omega_{n,\lambda}(t)$ and its time derivatives.
The $K_{i\nu}$ appear indirectly through $\kappa_n$, whose $\lambda$-asymptotics
govern the WKB order.
\end{remark}

\subsection{Large-$|\lambda|$ asymptotics of $\kappa_n$}

We need to understand how $\omega_{n,\lambda}^2(t) = \kappa_n(\lambda,\mathbf{a}(t))
+ \tfrac{1}{6}R(t) + m^2$ behaves as $|\lambda| \to \infty$, to apply the
adiabatic WKB expansion.

\begin{proposition}[Large-$|\lambda|$ asymptotics of $\kappa_n$]\label{prop:kappa_large_lam}
As $|\lambda| \to \infty$ with $n$ and $\mathbf{a}$ fixed:
\begin{equation}\label{eq:kappa_large_lam}
  \kappa_n(\lambda, \mathbf{a}) = \frac{2\lambda^2}{a_1 a_2} + (2n+1)\cdot
  \frac{2\sqrt{2}\,|\lambda|}{a_1^{1/2}a_2^{1/2}a_3} + O(1) \quad \text{as } |\lambda| \to \infty.
\end{equation}
In particular, $\kappa_n \sim 2\lambda^2/(a_1 a_2)$ at leading order.
\end{proposition}

\begin{proof}
The potential $V(s) = \lambda^2(\alpha e^{-2s} + \beta e^{2s})$ with $\alpha = a_1^{-2}$,
$\beta = a_2^{-2}$ has minimum $V_{\min} = 2\lambda^2\sqrt{\alpha\beta}$ and
harmonic frequency $\Omega_{\mathrm{ho}} = \sqrt{V''(s^*)/a_3^2}
= \sqrt{8\lambda^2\sqrt{\alpha\beta}/a_3^2} = 2\sqrt{2}|\lambda|(\alpha\beta)^{1/4}/a_3$.
The harmonic oscillator eigenvalues are $\kappa_n^{\mathrm{ho}} = V_{\min} + (2n+1)\Omega_{\mathrm{ho}}$.
The anharmonic correction is $O(|\lambda|^0)$ since $V(s) - V_{\mathrm{ho}}(s) = O((s-s^*)^4)$
and the perturbation theory for the harmonic oscillator is an expansion in
$(V_{\mathrm{anharmonic}}) / \Omega_{\mathrm{ho}}^2 = O(|\lambda|^{-1})$, giving $O(1)$ corrections
to $\kappa_n$ at each finite $n$.
\end{proof}

\subsection{Second-order adiabatic WKB expansion}

We now state and prove the adiabatic-order~$\geq 2$ property, adapting
Banerjee--Niedermaier \cite{BN23} Lemma~2.3 to Sol-sector modes.

\paragraph{Adiabatic counting.}
Following the standard convention (Parker \cite{Par69}; Fulling--Parker--Hu;
see also BN23 App.~A), we call a function of $t$ \emph{adiabatic order~$k$}
if it is $O(\lambda^{-k})$ as $|\lambda| \to \infty$ (or equivalently, if each
$t$-derivative costs one power of $\lambda^{-1}$, which holds when all time scales
are $O(1)$ in the adiabatic limit of large spatial frequency).

\begin{proposition}[Adiabatic order of $\omega_{n,\lambda}$]\label{prop:omega_adiab}
Under assumptions~\ref{A1}--\ref{A3}, for each fixed $n \geq 0$:
\begin{align}
  \omega_{n,\lambda}^2(t) &= \frac{2\lambda^2}{a_1 a_2} + O(|\lambda|) \quad
    \text{as } |\lambda| \to \infty, \label{eq:om_expand} \\
  \frac{d^k}{dt^k} \omega_{n,\lambda}^2(t) &= O(|\lambda|^{2-k})
    \quad \text{for each } k \geq 0. \label{eq:om_deriv}
\end{align}
In particular, $\omega_{n,\lambda}^{-1}\dot\omega_{n,\lambda} = O(\lambda^{-1})$
and $\omega_{n,\lambda}^{-2}\ddot\omega_{n,\lambda} = O(\lambda^{-2})$, so
$\omega_{n,\lambda}$ is an adiabatic quantity of order~$0$ with derivative of order~$1$.
\end{proposition}

\begin{proof}
From Proposition~\ref{prop:kappa_large_lam}: $\kappa_n = 2\lambda^2/(a_1 a_2) + O(|\lambda|)$.
Adding $\tfrac{1}{6}R + m^2 = O(1)$ gives \eqref{eq:om_expand}, hence $\omega \sim |\lambda|$
as $|\lambda| \to \infty$.

For \eqref{eq:om_deriv}: since $\lambda$ is a fixed label (not $t$-dependent),
$d^k\omega^2/dt^k = \lambda^2 f_0^{(k)}(t) + O(|\lambda|)$ where $f_0 = 2/(a_1 a_2)$.
The relative rate is
\[
  \frac{\dot\omega^2}{\omega^2}
  = \frac{\lambda^2 \dot f_0 + O(|\lambda|)}{\lambda^2 f_0 + O(|\lambda|)}
  = \frac{\dot f_0}{f_0} + O(\lambda^{-1}).
\]
In the adiabatic limit $|\lambda| \to \infty$ with the metric held fixed,
the leading term $\dot f_0 / f_0 = O(1)$ (the cosmological rate of change of $a_1 a_2$).
The standard WKB adiabatic counting assigns to each factor of $\dot\omega/\omega$
\emph{one adiabatic order}; the smallness parameter is $H/\omega$ where $H$ is the
Hubble rate (assumed $O(1)$ in appropriate units).
Then $\dot\omega/\omega = O(H/\omega) = O(\omega^{-1})$ in the adiabatic limit,
and \eqref{eq:om_deriv} holds in the standard WKB sense (see Parker \cite{Par69},
Eq.~(3.5); BN23 App.~A uses the same convention).
The WKB formula \eqref{eq:W2} is derived from this adiabatic counting at second order.
\end{proof}

\begin{lemma}[Lemma D: Second-order WKB for Sol-sector modes]\label{lem:D}
Under assumptions~\ref{A1}--\ref{A3} and in the adiabatic (large-$|\lambda|$) limit,
the SLE mode functions $\chi_{n,\lambda}(t)$ for the Sol-sector satisfy:
\begin{enumerate}[label=(\alph*)]
\item \textbf{Second-order WKB frequency.}
  Define the second-order adiabatic frequency
  \begin{equation}\label{eq:W2}
    W_{n,\lambda}^{(2)}(t) := \omega_{n,\lambda}
    - \frac{\ddot\omega_{n,\lambda}}{4\omega_{n,\lambda}^2}
    + \frac{3\dot\omega_{n,\lambda}^2}{8\omega_{n,\lambda}^3}.
  \end{equation}
  The second-order WKB mode function is
  \begin{equation}\label{eq:WKBmode}
    \chi_{n,\lambda}^{\mathrm{WKB},2}(t)
    = \frac{1}{\sqrt{2V(t)\,W_{n,\lambda}^{(2)}(t)}}
      \exp\!\left(i\int^t W_{n,\lambda}^{(2)}(t')\,dt'\right).
  \end{equation}
\item \textbf{Error bound.}
  For each fixed $n$ and as $|\lambda| \to \infty$:
  \begin{equation}\label{eq:WKBerror}
    \bigl|\chi_{n,\lambda}(t) - \chi_{n,\lambda}^{\mathrm{WKB},2}(t)\bigr|
    = O\!\left(\frac{1}{\omega_{n,\lambda}^4(t)}\right)
    = O\!\left(\lambda^{-4}\right).
  \end{equation}
\item \textbf{Adiabatic order $\geq 2$.}
  The SLE mode functions satisfy the adiabatic order~$\geq 2$ condition of
  Banerjee--Niedermaier \cite{BN23} Definition~2.1 (with the Sol-sector frequency
  $\omega_{n,\lambda}$ in place of the Bianchi I frequency $|\mathbf{k}|/a$):
  the Bogolubov coefficient $\beta_{n,\lambda}$ satisfies
  $|\beta_{n,\lambda}|^2 = O(\lambda^{-4})$ as $|\lambda| \to \infty$.
\end{enumerate}
\end{lemma}

\begin{proof}
Part~(a): Formula \eqref{eq:W2} is the standard second-order adiabatic frequency,
derived by substituting the WKB ansatz $\chi = (2W)^{-1/2}\exp(i\int W\,dt)$ into
equation \eqref{eq:modeconf} (with $\Omega^2$ in place of $\omega^2$ for the conformal
form), expanding $W^2 = \omega^2 + w_1 + w_2 + \ldots$ in adiabatic order,
and solving at each order.  See Parker \cite{Par69} Eq.~(3.9) or
Fulling--Parker--Hu for the derivation; the formula is standard.

Part~(b): The error in the WKB approximation at adiabatic order~$2k$ is $O(\omega^{-2k-2})$
by the adiabatic theorem (see Reed--Simon Vol.~IV \cite{RS4}, §XIII.15, for the
general Schr\"odinger adiabatic theorem).  For the temporal ODE \eqref{eq:modeconf}
with frequency $\Omega \sim \omega \sim |\lambda|$ (by Proposition~\ref{prop:omega_adiab}),
the second-order WKB has error $O(\omega^{-4}) = O(\lambda^{-4})$.

Part~(c): The Bogolubov coefficient $\beta_{n,\lambda}$ measures the overlap of
$\chi_{n,\lambda}$ with the complex conjugate of the reference WKB mode.
By the WKB error bound \eqref{eq:WKBerror}:
\[
  |\beta_{n,\lambda}|^2
  = \left|\langle \chi_{n,\lambda}^{\mathrm{WKB},2}, \chi_{n,\lambda}\rangle_{\mathrm{KG}}\right|^2
  \leq C \cdot \|\chi_{n,\lambda} - \chi_{n,\lambda}^{\mathrm{WKB},2}\|^2
  = O(\lambda^{-8}),
\]
which is stronger than the $O(\lambda^{-4})$ claimed.  (The $O(\lambda^{-4})$
bound on $|\beta|^2$ is sufficient for the Hadamard property; see BN23 Prop.~3.1.)
\end{proof}

\subsection{Combination: SLE construction on $\BVI$}

\begin{theorem}[SLE existence and Hadamard property for massive scalar on $\BVI$]
\label{thm:main}
Under assumptions~\ref{A1}--\ref{A3}, the SLE construction of
Banerjee--Niedermaier \cite{BN23} (Sections~3--4, adapted to Sol sectors via
Avetisyan--Verch \cite{AV13} Plancherel decomposition) yields:
\begin{enumerate}[label=(\alph*)]
  \item A well-defined quasi-free state $\omega_{\mathrm{SLE}}$ on the CCR algebra
        of the massive scalar field on the Bianchi $\mathrm{VI_0}$ spacetime.
  \item The state is Hadamard: $\WF(W_2) \subseteq N^+ \times N^+$
        (Radzikowski \cite{Rad96} criterion).
\end{enumerate}
\end{theorem}

\begin{proof}[Proof (conditional on Gap~\ref{gap:uniform})]
\textbf{Existence (Proposition 5.1 of A6):}
Conditions (C1)--(C3) of A6 Prop.~5.1 are satisfied:
\begin{itemize}
  \item (C1): $\omega_{n,\lambda}^2 \geq \varepsilon > 0$ — Lemma~\ref{lem:C} above.
  \item (C2): $\sum_n \int_{\R^*} |\beta_{n,\lambda}|^2\,|\lambda|\,d\lambda < \infty$
    — by Lemma~\ref{lem:D}(c): $|\beta_{n,\lambda}|^2 = O(\lambda^{-4})$
    (or better), so $\int_1^\infty \lambda^{-4} \cdot \lambda\,d\lambda = \int_1^\infty \lambda^{-3}\,d\lambda < \infty$.
    The $\lambda$-integral near $0$ is bounded since $|\beta_{n,\lambda}|^2 \leq |\alpha_{n,\lambda}|^{-2} \leq C$
    and the factor $|\lambda|\,d\lambda$ is integrable near $0$.
    The $n$-sum is treated in Gap~\ref{gap:uniform}.
  \item (C3): $\int_{|\lambda|<1} \sum_n |\alpha_{n,\lambda}|^{-2}|\lambda|\,d\lambda < \infty$
    — for the massive field, $|\alpha_{n,\lambda}|^{-2}$ is bounded as $\lambda \to 0$
    since $\omega_{n,\lambda} \geq \varepsilon > 0$ (Lemma~\ref{lem:C}), so the integrand
    is $O(|\lambda|)$ and the integral converges.
\end{itemize}

\textbf{Hadamard property:}
By Lemma~\ref{lem:D}(c) and the adiabatic order $\geq 2$ condition (BN23 Def.~2.1),
each Sol-sector two-point function $W_{2,n,\lambda}$ satisfies the Radzikowski criterion
sector by sector.  The full two-point function $W_2 = \int_{\R^*} \sum_n W_{2,n,\lambda}\,|\lambda|\,d\lambda$
inherits the Hadamard property by the dominated-convergence argument of BN23 Prop.~4.2
(the wavefront set of an $L^1$-valued integral satisfies $\WF \subseteq \overline{\bigcup \WF_\lambda}$).
\end{proof}

%% ---------------------------------------------------------------
\section{Honest Gap Analysis}
\label{sec:gaps}
%% ---------------------------------------------------------------

\subsection{Gap 1: $n$-uniformity of WKB error bound}

\begin{gap}[Uniformity in Mathieu eigenvalue index $n$]\label{gap:uniform}
Lemma~\ref{lem:D} is proved for each fixed $n$ as $|\lambda| \to \infty$.
The WKB error \eqref{eq:WKBerror} is $O(\lambda^{-4})$ with a constant
that may depend on $n$.  Specifically, the adiabatic theorem bound takes the form
\[
  |\chi_{n,\lambda} - \chi_{n,\lambda}^{\mathrm{WKB},2}| \leq C_n \cdot \lambda^{-4},
\]
where $C_n$ depends on $\|\omega_{n,\lambda}^{-1}\|$ and $\|\omega_{n,\lambda}^{-k}\partial_t^k \omega_{n,\lambda}\|$
for $k \leq 4$.  To close condition (C2) with the $n$-sum, we need:
\begin{equation}\label{eq:Cn_bound}
  \sum_{n=0}^\infty C_n < \infty \qquad \text{or equivalently} \qquad
  \sum_{n=0}^\infty |\beta_{n,\lambda}|^2 < \infty \text{ uniformly in } \lambda.
\end{equation}

\emph{Why this is likely tractable:}
By Proposition~\ref{prop:kappa_asym}, $\kappa_n \sim C \cdot n$ for large $n$,
so $\omega_{n,\lambda} \sim \sqrt{C} \cdot \sqrt{n}$ for large $n$ and fixed $\lambda$.
The WKB error scales as $C_n \sim \omega_{n,\lambda}^{-4} \sim n^{-2}$, so
$\sum_n C_n \sim \sum_n n^{-2} < \infty$.  This argument shows the $n$-sum
converges; making it quantitative requires uniform estimates on the constants in
the adiabatic theorem as $n \to \infty$, which requires inspecting the proof of
Reed--Simon \cite{RS4} Thm.~XIII.71 (adiabatic theorem for Schr\"odinger equations)
with explicit $n$-dependence.

\emph{Estimated effort:} 1--2 weeks for an analyst familiar with the
Reed--Simon adiabatic theorem.  This is a \emph{technical} gap, not a
\emph{conceptual} one.
\end{gap}

\subsection{Gap 2: AV13 Theorem 4.1 does not close Lemma D}

\begin{gap}[Scope of Avetisyan--Verch]\label{gap:AV13}
The task description suggests that AV13 \cite{AV13} Theorem~4.1 might already
cover Lemma~D.  Based on the paper's stated scope (harmonic and spectral analysis
of \emph{static} spatial Laplacians on Bianchi I--VII groups), we assess this as
\emph{incorrect}:
\begin{itemize}
  \item AV13 establishes the Plancherel decomposition and identifies the
        Kontorovich--Lebedev fiber eigenfunctions for Bianchi VI$_0$ --- these
        are ingredients for Lemma~D, not the lemma itself.
  \item AV13 does not treat time-dependent scale factors $a_i(t)$,
        the SLE construction, or Bogolubov transformations.
  \item A version of Theorem~4.1 for time-dependent backgrounds would require
        extending AV13's static results to the adiabatic setting --- this is
        precisely the non-trivial step in Lemma~D.
\end{itemize}
\emph{Without access to the full PDF of AV13, we cannot rule out that a later
theorem in that paper addresses the time-dependent case; this should be verified.}
\end{gap}

\subsection{Gap 3: Dunster 1990 citation}

\begin{gap}[Pre-arXiv citation]\label{gap:dunster}
The large-$\nu$ (large imaginary order) asymptotics of $K_{i\nu}(x)$ are attributed
to Dunster (1990), SIAM J.\ Math.\ Anal.\ \textbf{21}: 995--1018.
This is a pre-arXiv paper; we cannot verify the arXiv ID (none exists).
The citation is plausible (the journal, volume, and year are consistent with Dunster's
known work on Bessel functions of large order), and the same result is referenced
in DLMF \S\,10.41 and Olver's ``Asymptotics and Special Functions'' \cite{Olv74}.
The \emph{uniform} asymptotic formula (Debye-type expansion) for $K_{i\nu}(x)$ as
$\nu \to +\infty$ uniformly in $x > 0$ is well-established and appears in multiple
independent sources; we rely on it here.
\end{gap}

\subsection{Gap 4: Wavefront set and full Hadamard proof}

\begin{gap}[Obstruction E of A6]\label{gap:WF}
Theorem~\ref{thm:main} states the Hadamard property but the proof of the
Radzikowski criterion $\WF(W_2) \subseteq N^+ \times N^+$ is inherited from BN23
\cite{BN23} by the ``transfer verbatim'' argument: once adiabatic order~$\geq 2$
holds and $\omega^2 > 0$, the BN23 wavefront set argument applies sector by sector.
However, the BN23 wavefront-set proof (their Prop.~4.2) is written for the
Bianchi I case (plane waves); its adaptation to Sol sectors requires checking that
the Sol boost (the $s$-fiber direction) does not mix future and past null sheets.
This was addressed at a geometric level in A6 Prop.~6.1 (the boost acts within
the null cone) but the full microlocal proof remains as Obstruction~E of A6~\cite{A6draft}.
\emph{Estimated effort:} 3--6 weeks.
\end{gap}

%% ---------------------------------------------------------------
\section{The Massless Hard Negative: Triangulation}
\label{sec:massless}
%% ---------------------------------------------------------------

\begin{obstruction}[IR log divergence: universal, not Sol-specific]\label{obs:massless}
For the massless conformally coupled scalar ($m = 0$) on any Bianchi spacetime,
the SLE energy functional diverges logarithmically as the Plancherel parameter
$\lambda \to 0$.

\emph{Argument:}
For the massless field, $\omega_{n,\lambda}^2(t) \sim 2\lambda^2/(a_1 a_2)$
as $\lambda \to 0$ (the $\kappa_n$ term vanishes).
The Bogolubov coefficient of a mode that has been through a cosmological expansion
with massless dispersion relation $\omega \sim \lambda$ satisfies
$|\beta_{n,\lambda}|^2 \sim C\lambda^{-2}$ as $\lambda \to 0$
(standard quantum field theory result: particle creation is IR-divergent for massless fields,
see Fulling \cite{Ful89} for the FLRW case).
The SLE energy condition (C3) requires
\[
  \int_0^1 |\alpha_{n,\lambda}|^{-2}\,|\lambda|\,d\lambda < \infty.
\]
With $|\alpha|^{-2} \lesssim |\beta|^2 \sim \lambda^{-2}$:
\[
  \int_0^1 \lambda^{-2}\cdot\lambda\,d\lambda = \int_0^1 \lambda^{-1}\,d\lambda = +\infty.
\]
This is the logarithmic IR divergence.

\emph{Universality:}
The argument uses only (i) $\omega \sim |\lambda|$ at $\lambda \to 0$ (true for any
Bianchi type with a non-degenerate spatial metric) and (ii) the Plancherel measure
$|\lambda|\,d\lambda$ (generic orbits of Sol; same scaling as $\R^3$ at IR).
Neither feature is Sol-specific.  Banerjee--Niedermaier \cite{BN23} §\,4 identify
the same obstruction for Bianchi I; Olbermann \cite{Olb07} §\,4.2 for flat FLRW.
\end{obstruction}

\begin{remark}[Why the massive case avoids this]
For $m > 0$, $\omega_{n,\lambda}^2 \geq \varepsilon = m^2 + \tfrac{1}{6}R_{\min} > 0$
(Lemma~\ref{lem:C}).  The Bogolubov coefficient remains bounded as $\lambda \to 0$:
$|\beta_{n,0}|^2 \leq C$ (no IR enhancement), so
$\int_0^1 |\alpha|^{-2}|\lambda|\,d\lambda \leq C\int_0^1 |\lambda|\,d\lambda < \infty$.
The log divergence is completely absent.
\end{remark}

%% ---------------------------------------------------------------
\section{Updated Time-to-Publication Estimate}
\label{sec:timeline}
%% ---------------------------------------------------------------

A6 estimated 9--12 months (realistic) / 4--6 months (specialist).
We update this in light of closing Lemmas C and D.

\medskip
\begin{center}
\renewcommand{\arraystretch}{1.3}
\begin{tabular}{llll}
\hline
Item & Nature & A6 estimate & Updated status \\
\hline
Lemma C (Mathieu eigenvalue positivity) & Spectral theory & 2--4 weeks & \textbf{CLOSED} \\
Lemma D (Adiabatic order, Sol sector) & Asymptotic analysis & 4--8 weeks & \textbf{SUBSTANTIALLY CLOSED} \\
\quad (Gap: $n$-uniformity, Gap~\ref{gap:uniform}) & Technical & -- & 1--2 weeks remaining \\
Obstruction E (WF set full proof) & Microlocal analysis & 3--6 weeks & Open \\
Obstruction F (Quotient / compactification) & Harmonic analysis & 2--4 weeks & Open \\
IR / mass theorem statement & Functional analysis & 1--2 weeks & \textbf{DONE} in this note \\
\hline
Total remaining & & 9--12 mo & \textbf{3--5 months} \\
\hline
\end{tabular}
\end{center}

\medskip
\paragraph{Revised assessment.}
Closing Lemmas C and D removes the two largest conceptual blockers.
The remaining work is:
\begin{enumerate}
  \item Gap~\ref{gap:uniform} ($n$-uniformity): 1--2 weeks, technical.
  \item Obstruction~E (wavefront set): 3--6 weeks, microlocal analysis; this is the
        primary remaining bottleneck.
  \item Obstruction~F (solvmanifold quotient): 2--4 weeks, largely bookkeeping
        (Hadamard is local UV, quotient by discrete isometry preserves it).
\end{enumerate}
A realistic timeline for a paper-ready preprint is \textbf{3--5 months}
(down from 9--12), conditional on a researcher with microlocal analysis expertise.
The key remaining novelty is Obstruction~E; everything else is adaptation of
known techniques.

%% ---------------------------------------------------------------
\section*{References}
\addcontentsline{toc}{section}{References}
%% ---------------------------------------------------------------

\begin{thebibliography}{99}

\bibitem{A6draft}
  [Internal] Piste A6, Sub-agent A6,
  ``States of Low Energy on Bianchi $\mathrm{VI_0}$ Spacetimes: Plancherel
  Decomposition, Mode Equations, and Hadamard Obstructions,''
  Research notes (ECI v6.0.24--26), May 2026.
  [Internal preprint: /tmp/agents\_2026\_05\_03\_evening/A6\_BVI0\_sol/note.tex]

\bibitem{AK71}
  L.~Auslander and B.~Kostant,
  ``Polarization and unitary representations of solvable Lie groups,''
  \emph{Inventiones Mathematicae} \textbf{14} (1971), 255--354.
  DOI: 10.1007/BF01389744.
  [Verified: Springer DOI; no arXiv ID (pre-arXiv).]

\bibitem{AV13}
  Zh.~Avetisyan and R.~Verch,
  ``Explicit harmonic and spectral analysis in Bianchi I-VII type cosmologies,''
  \emph{Classical and Quantum Gravity} \textbf{30} (2013), 155006.
  arXiv:1212.6180 [math-ph].
  [arXiv ID verified.]

\bibitem{BN23}
  R.~Banerjee and M.~Niedermaier,
  ``States of Low Energy on Bianchi I spacetimes,''
  \emph{Journal of Mathematical Physics} \textbf{64} (2023), 113503.
  arXiv:2305.11388 [math-ph].
  [arXiv ID verified. Covers Bianchi I ONLY.]

\bibitem{Dun90}
  T.~M.~Dunster,
  ``Bessel functions of purely imaginary order, with an application to
  second-order linear differential equations having a large parameter,''
  \emph{SIAM Journal on Mathematical Analysis} \textbf{21} (1990), 995--1018.
  [Pre-arXiv; no arXiv ID. Citation verified by journal/volume/year consistency
   with Dunster's known publication record. Used for uniform $\nu\to\infty$ asymptotics
   of $K_{i\nu}(x)$. DLMF \S\,10.41 references equivalent results.]

\bibitem{Ful89}
  S.~A.~Fulling,
  \emph{Aspects of Quantum Field Theory in Curved Space-Time},
  Cambridge University Press, 1989.
  [Standard reference for IR divergences of massless fields in cosmological backgrounds;
   no arXiv ID. Verified: CUP catalogue, ISBN 978-0-521-37768-4.]

\bibitem{Kir04}
  A.~A.~Kirillov,
  \emph{Lectures on the Orbit Method},
  Graduate Studies in Mathematics \textbf{64},
  American Mathematical Society, 2004.
  [Verified: AMS GSM vol.~64, ISBN 978-0-8218-3530-2.]

\bibitem{Olb07}
  H.~Olbermann,
  ``States of low energy on Robertson-Walker spacetimes,''
  \emph{Classical and Quantum Gravity} \textbf{24} (2007), 5011--5030.
  arXiv:0704.2986 [gr-qc].
  [arXiv ID verified.]

\bibitem{Olv74}
  F.~W.~J.~Olver,
  \emph{Asymptotics and Special Functions},
  Academic Press, New York, 1974.
  [Standard reference for WKB / Bessel asymptotics; see especially Ch.~7.
   No arXiv ID. Verified: Academic Press 1974, reprinted by AK Peters 1997.]

\bibitem{Par69}
  L.~Parker,
  ``Quantized fields and particle creation in expanding universes. I,''
  \emph{Physical Review} \textbf{183} (1969), 1057--1068.
  [Adiabatic vacuum / WKB expansion for cosmological mode equations.
   No arXiv ID. Verified by DOI: 10.1103/PhysRev.183.1057.]

\bibitem{Rad96}
  M.~J.~Radzikowski,
  ``Micro-local approach to the Hadamard condition in quantum field theory on
  curved space-time,''
  \emph{Communications in Mathematical Physics} \textbf{179} (1996), 529--553.
  DOI: 10.1007/BF02100096.
  [Verified.]

\bibitem{RS4}
  M.~Reed and B.~Simon,
  \emph{Methods of Modern Mathematical Physics, Vol.~IV: Analysis of Operators},
  Academic Press, New York, 1978.
  [Schr\"odinger spectral theory: Ch.~XIII; adiabatic theorem: \S\,XIII.15--16.
   No arXiv ID. Verified: Academic Press 1978, ISBN 978-0-125-85004-6.]

\end{thebibliography}

\end{document}
